\clearpage
\VerdierPhysicalPageAnchor{0262}%
\begingroup
\small
\setlength{\parindent}{0pt}%
\setlength{\parskip}{0pt}%
\setlength{\emergencystretch}{2em}%
\newcommand{\VerdierTermRow}[3]{%
  \VerdierSemanticAnchor{#1}%
  \par\noindent\hangindent=1.5em\hangafter=1
  #2。\enspace\textup{#3。}\par\smallskip}

\VerdierTermRow{idxt-p262-r01}{$\mathcal{F}\ell_n(\widetilde{\mathbb Z})$ 的有限对象}{第二章，4.3.1}
\VerdierTermRow{idxt-p262-r02}{$\operatorname{comp}^{n}(C)$ 的可求和对象}{第一章，2.1.7}
\VerdierTermRow{idxt-p262-r03}{$\displaystyle\prod_{i\in[n]}\operatorname{comp}(C_i)$ 的可求和对象}{第一章，2.4.2}
\VerdierTermRow{idxt-p262-r04}{谱对象}{第二章，4.3.2}
\VerdierTermRow{idxt-p262-r05}{取值于 $D(\mathcal A)$ 的谱对象，由按第 $m$ 次滤过的复形定义}{第三章，4.1.8}
\VerdierTermRow{idxt-p262-r06}{取值于 $K(\mathcal A)$ 的谱对象，由按第 $m$ 次滤过的复形定义}{第三章，4.1.7}
\VerdierTermRow{idxt-p262-r07}{与取值于加法范畴复形范畴的函子关联的谱对象}{第二章，4.1.3}
\VerdierTermRow{idxt-p262-r08}{Cartan--Eilenberg 谱对象}{第三章，4.6.7}
\VerdierTermRow{idxt-p262-r09}{取值于阿贝尔范畴的 $\mathcal J$ 型谱对象}{第二章，4.1.4}
\VerdierTermRow{idxt-p262-r10}{取值于三角范畴的 $\mathcal J$ 型谱对象}{第二章，4.1.2}
\VerdierTermRow{idxt-p262-r11}{定常谱对象}{第二章，4.4.2}
\VerdierTermRow{idxt-p262-r12}{符号群在取值于加法范畴的函子上的典范作用}{第一章，1.4.3}
\VerdierTermRow{idxt-p262-r13}{$\Pi$-扩张（函子 $\int_\varphi$）}{第一章，2.2.3}
\VerdierTermRow{idxt-p262-r14}{函子到复形同伦范畴的 $\Pi$-扩张}{第一章，2.5.10}
\VerdierTermRow{idxt-p262-r15}{多加法函子的 $\Pi$-扩张}{第一章，2.4.2}
\VerdierTermRow{idxt-p262-r16}{函子态射到复形同伦范畴的 $\Pi$-扩张}{第一章，2.5.10}
\VerdierTermRow{idxt-p262-r17}{复形所关联的取值于 $D(\mathcal A)$ 的第一谱对象}{第三章，4.1.5}
\VerdierTermRow{idxt-p262-r18}{复形所关联的取值于 $K(\mathcal A)$ 的第一谱对象}{第三章，4.1.4}
\VerdierTermRow{idxt-p262-r19}{与复形关联的第一谱对象}{第三章，4.1.5}
\VerdierTermRow{idxt-p262-r20}{取值于 $D(\mathcal A)$ 的第一典范谱对象}{第三章，4.1.5}
\VerdierTermRow{idxt-p262-r21}{取值于 $K(\mathcal A)$ 的第一典范谱对象}{第三章，4.1.4}
\endgroup
